\section*{Capítulo \ref{ch:b2:chromatography} --- Cromatografía}

\begin{solution}{exo:b2:chromatography:1}
$k = (5.0 - 1.0)/1.0 = 4.0$; fracción de tiempo en la \omterm{def:b2:chromatography:principle}{fase móvil} $1/(1 + k) = 0.20$.
\end{solution}

\begin{solution}{exo:b2:chromatography:2}
$N = 16(6.0/0.30)^2 = 16 \times 400 = 6400$.
\end{solution}

\begin{solution}{exo:b2:chromatography:3}
$H = \qty{250}{mm}/10\,000 = \qty{0.025}{mm} = \qty{25}{\micro m}$.
\end{solution}

\begin{solution}{exo:b2:chromatography:4}
Etanol en sangre: CG (volátil). Proteína: HPLC. Cafeína en una bebida: HPLC (en agua, no lo bastante volátil
sin preparación). Benceno en gasolina: CG. Azúcares: HPLC (se descomponen antes de hervir).
\end{solution}

\begin{solution}{exo:b2:chromatography:5}
$R_s = 2(4.6 - 4.0)/(0.40 + 0.50) = 1.2/0.90 \approx 1.3$: no están del todo separados hasta la línea de base (menos de 1.5).
\end{solution}

\begin{solution}{exo:b2:chromatography:6}
$\sqrt N = 4R_s\,\dfrac{\alpha}{\alpha - 1}\,\dfrac{1 + k_2}{k_2} = 4 \times 1.5 \times 11 \times 1.25 =
82.5$, luego $N \approx 6.8 \times 10^3$.
\end{solution}

\begin{solution}{exo:b2:chromatography:7}
$u_{\mathrm{opt}} = \sqrt{8/2} = \qty{2}{mm/s}$; $H_{\min} = 5 + 2\sqrt{16} = \qty{13}{\micro m}$.
\end{solution}

\begin{solution}{exo:b2:chromatography:8}
Fenol (\ce{OH} polar, enlaces de hidrógeno con el agua) primero, luego benceno, luego tolueno (un \ce{CH3}
más, más apolar). Más metanol hace la \omterm{def:b2:chromatography:principle}{fase móvil} menos polar: los tres eluyen antes, en
el mismo orden.
\end{solution}

\begin{solution}{exo:b2:chromatography:9}
$F = (860/400)/(20.0/10.0) = 2.15/2.00 = 1.075$. Muestra: $c_X = (645/410) \times 10.0/1.075 \approx
\qty{14.6}{mg/L}$.
\end{solution}

\begin{solution}{exo:b2:chromatography:10}
$R_s \propto \sqrt N \propto \sqrt L$: hay que multiplicar $L$ por $(1.5/1.06)^2 \approx 2.0$, lo que duplica
el tiempo de análisis y la presión. Otras palancas: la selectividad (\omterm{def:b2:chromatography:principle}{fase móvil}, \omterm{def:b2:chromatography:principle}{fase estacionaria},
temperatura), partículas más pequeñas, el caudal óptimo.
\end{solution}

\begin{solution}{exo:b2:chromatography:11}
Los picos están separados $4\sigma R_s$, de modo que el punto medio está a $2\sigma R_s$ de cada uno. $R_s = 1$: cada pico
contribuye con $\mathrm e^{-(2)^2/2} = \mathrm e^{-2} = 0.135$; el valle está a $0.27$ de la altura de pico.
$R_s = 1.5$: $2\mathrm e^{-(3)^2/2} = 2\mathrm e^{-4.5} \approx 0.022$, alrededor del 2~\%: de vuelta a la línea de base
a efectos prácticos.
\end{solution}

\begin{solution}{exo:b2:chromatography:12}
$k_2/(1+k_2)$: 0.67, 0.83, 0.91; tiempo en unidades de $t_M$: 3, 6, 11. Pasar de 2 a 5 gana un 25~\% de
\omterm{def:b2:chromatography:separation}{resolución} por el doble de tiempo; de 5 a 10, un 9~\% por casi el doble otra vez. Un $k$ de 2 a 5, aproximadamente, es
el compromiso habitual.
\end{solution}

\begin{solution}{pb:b2:chromatography:1}
\textbf{1.} Estacionaria: granos de sílice que llevan largas cadenas alquílicas, apolar. Móvil: agua con
metanol, polar.
\textbf{2.} Son muy polares y permanecen en la \omterm{def:b2:chromatography:principle}{fase móvil}: $k \approx 0$.
\textbf{3.} La cafeína tiene un grupo metilo más que la teofilina, y la teofilina, un \ce{N-H} que
forma enlaces de hidrógeno con el agua: la teofilina es más polar y eluye primero.
\textbf{4.} Es de estructura próxima (respuesta y comportamiento similares), está ausente de la bebida y eluye
cerca de la cafeína sin dejar de estar resuelta respecto a ella.
\textbf{5.} Ambos disminuirían: una \omterm{def:b2:chromatography:principle}{fase móvil} menos polar compite mejor con la \omterm{def:b2:chromatography:principle}{fase estacionaria}.
\textbf{6.} El tiempo que tarda la \omterm{def:b2:chromatography:principle}{fase móvil} en atravesar la columna, el tiempo que un compuesto pasa avanzando.
\textbf{7.} $k_{\mathrm{theo}} = (3.0 - 1.0)/1.0 = 2.0$; $k_{\mathrm{caf}} = 2.4$.
\textbf{8.} $\alpha = 2.4/2.0 = 1.2$.
\textbf{9.} Cafeína: $N = 16(3.4/0.20)^2 = 4624 \approx 4.6 \times 10^3$; teofilina: $16(3.0/0.18)^2
\approx 4.4 \times 10^3$.
\textbf{10.} $H = \qty{150}{mm}/4624 \approx \qty{0.032}{mm} = \qty{32}{\micro m}$.
\textbf{11.} $R_s = 2(3.4 - 3.0)/(0.18 + 0.20) = 0.80/0.38 \approx 2.1$.
\textbf{12.} $\frac{\sqrt{4624}}{4} \times \frac{0.2}{1.2} \times \frac{2.4}{3.4} = 17 \times 0.167 \times
0.706 \approx 2.0$; la pequeña diferencia procede de la hipótesis de anchuras iguales de la ecuación.
\textbf{13.} Sí: $R_s > 1.5$.
\textbf{14.} $N$ se reduce a la mitad: $R_s \approx 2.1/\sqrt2 \approx 1.5$, justo en la separación hasta la línea de base.
\textbf{15.} Se reduce a la mitad: la cafeína sale a \qty{1.7}{min}.
\textbf{16.} $R_s \approx 2.1\sqrt2 \approx 3.0$, por el doble de tiempo y el doble de presión: aquí es un despilfarro.
\textbf{17.} La composición de la \omterm{def:b2:chromatography:principle}{fase móvil} (fracción de metanol, pH, otro disolvente orgánico) o la
\omterm{def:b2:chromatography:principle}{fase estacionaria}.
\textbf{18.} Poco: $k_2/(1 + k_2) = 0.71$ ya; con $k_2 = 5$ sería 0.83 (+18~\%) para un
tiempo de análisis multiplicado por $6/3.4 \approx 1.8$.
\textbf{19.} $F = (1500/1000)/(50.0/40.0) = 1.50/1.25 = 1.20$.
\textbf{20.} Ambos picos proceden de la misma inyección: el volumen se cancela en la razón de áreas.
\textbf{21.} $c = (970/1010) \times 40.0/1.20 \approx \qty{32.0}{mg/L}$.
\textbf{22.} Dilución al décimo: \qty{320}{mg/L}.
\textbf{23.} $A_{IS}$ sería demasiado grande, y el resultado de cafeína, demasiado bajo.
\textbf{24.} Una serie de patrones de cafeína inyectados exactamente en las mismas condiciones que la muestra, una
recta de calibrado del área frente a la concentración, y volúmenes de inyección reproducibles.
\textbf{25.} $\qty{320}{mg/L} \times \qty{0.250}{L}$: $\boldsymbol{\approx \qty{80}{mg}}$ de cafeína por lata (datos
de ejercicio).
data).
\end{solution}
